% answers_short.tex —— 1989 年·广东卷文科候选卷 · 简短答案速查
% 由 tools/make_short_answers.py 自 answers.tex 的 \gksolhead/\gkanswer 提取，
% 并与 1989/1989_数学真题_广东文_答案要点.md 逐题比对后生成，勿手改（改源 MD 后重跑）。

\begin{gknotice}
\textbf{使用说明}：本节只给最终答案便于快速核对，完整推导见解析版。答案为编写者独立推导，\textbf{非官方标准答案}。
\end{gknotice}

\begin{gknotice}
选择题答案中的 U 表示待核，未给出猜测答案，详情见解析版。
\end{gknotice}

\gkpart{选择题}
\begin{center}\begin{tabular}{cccccccccccccc}
\toprule
题号 & 1 & 2 & 3 & 4 & 5 & 6 & 7 & 8 & 9 & 10 & 11 & 12 & 13 \\
\midrule
答案 & A & C & D & C & A & B & D & D & U & A & B & B & A \\
\bottomrule
\end{tabular}\end{center}\medskip
\begin{center}\begin{tabular}{ccccccccccccc}
\toprule
题号 & 14 & 15 & 16 & 17 & 18 & 19 & 20 & 21 & 22 & 23 & 24 & 25 \\
\midrule
答案 & B & C & U & A & D & C & B & U & B & D & B & B \\
\bottomrule
\end{tabular}\end{center}

\gkpart{填空题}
\begin{itemize}[leftmargin=2.6em]
  \item[\textbf{一.}] (1) $-1$；(2) $y=\dfrac{1-x}{3x}$，$x\ne0$；(3) $3x+\sqrt3\,y=12$；(4) $\dfrac12$；(5) $23$
\end{itemize}

\gkpart{解答题}
\begin{itemize}[leftmargin=2.6em]
  \item[\textbf{二.}] (1) 见解析版；(2) $x=-\dfrac12$
  \item[\textbf{三.}] (1) $AC=BC$；(2) $SC=5\text{ cm}$
  \item[\textbf{四.}] (1) $-1-i=\sqrt2\left(\cos\dfrac{5\pi}{4}+i\sin\dfrac{5\pi}{4}\right)$；(2) $a\in\left(-1,-\dfrac12\right]\cup\left[-\dfrac15,+\infty\right)$
  \item[\textbf{五.}] (1) $x^2+\dfrac{(y+1)^2}{4}=1$；(2) $\theta=0$ 或 $\pi$；最大公共弦长为 $2$
\end{itemize}
